\documentclass[11pt, a4paper]{article}

\usepackage[a4paper, top=2cm, bottom=2cm, left=1.5cm, right=1.5cm]{geometry}
\usepackage{lmodern}
\usepackage{fontenc}
\usepackage{inputenc}
\usepackage{array}
\usepackage{tabularx}
\usepackage{longtable}
\usepackage{booktabs}
\usepackage{colortbl}
\usepackage{xcolor}
\usepackage{makecell}
\usepackage{fancyhdr}
\usepackage{setspace}
\usepackage{parskip}
% Colors
\definecolor{headerblue}{RGB}{255,255,255}
\definecolor{tablegray}{RGB}{220,220,220}
\definecolor{lightgray}{RGB}{240,240,240}

% Page style
\pagestyle{fancy}
\fancyhf{}
\fancyhead[L]{\small AIA Lab Practicals}
\fancyhead[R]{\small ET23BTIT133}
\renewcommand{\headrulewidth}{0.4pt}

\setlength{\parindent}{0pt}

\begin{document}

% -------------------------------------------------------
% Cover / Header block (matching the image)
% -------------------------------------------------------

\vspace*{0.8cm}

{\fontsize{22pt}{26pt}\selectfont\bfseries Artificial Intelligence and Applications}

\vspace{0.3cm}
{\large\color{gray} Practical Work and Lab Assignments}

\vspace{1cm}

\begin{tabular}{@{}p{0.48\textwidth} p{0.48\textwidth}@{}}
  \textbf{Name:} Sree Saicharan Vadapalli   & \quad \quad \quad \quad \textbf{Subject Code:} BTIT13601 \\[4pt]
  \textbf{Branch:} Information Technology   & \quad \quad \quad \quad \textbf{Year:} B. Tech III \textbf{(Sem VI)} \\[4pt]
  \textbf{Enrollment No:} ET23BTIT133       & \\
\end{tabular}

\vspace{1.5cm}

% -------------------------------------------------------
% INDEX heading
% -------------------------------------------------------

\begin{center}
  {\large\bfseries INDEX}
\end{center}

\vspace{0.5cm}

% -------------------------------------------------------
% Index Table
% -------------------------------------------------------

\renewcommand{\arraystretch}{1.6}

\begin{longtable}{
  |>{\centering\arraybackslash}p{0.7cm}
  |>{\centering\arraybackslash}p{1.7cm}
  |>{\arraybackslash}p{9.3cm}
  |>{\centering\arraybackslash}p{1.3cm}
  |>{\centering\arraybackslash}p{2.2cm}|
}

\hline
\rowcolor{tablegray}
\makecell{\textbf{Sr.}\\\textbf{No.}} & \textbf{Date} & \textbf{Experiments} & \makecell{\textbf{Page}\\\textbf{No.}} & \textbf{Signature} \\
\hline
\endfirsthead

\hline
\rowcolor{tablegray}
\makecell{\textbf{Sr.}\\\textbf{No.}} & \textbf{Date} & \textbf{Experiments} & \makecell{\textbf{Page}\\\textbf{No.}} & \textbf{Signature} \\
\hline
\endhead

1. & &
Study and make a detailed note on the following AI Projects:\newline
\quad i)   Deep Blue (chess program)\newline
\quad ii)  Chinook (checkers program)\newline
\quad iii) ALVINN (autonomous driving)\newline
\quad iv)  Eliza, ChatGPT (chatbots)\newline
\quad v)   Google Translate (machine translation)\newline
\quad vi)  Gemini (AI Assistant)\newline
\quad vii) Nano Banana (AI image generator)
& & \\
\hline

2. & &
Write a program to solve the 8-puzzle. First, check if the puzzle is solvable. If solvable, starting from any input state, it should display the sequence of moves to reach the ordered output state. Rule to move: Slide a tile from a non-empty square to an adjoining empty square.
& & \\
\hline

3. & &
Implement the WaterJug Problem. The state of the system is represented as $(x,y)$, where $x$ is the quantity of water in the 4-gallon jug and $y$ is the quantity of water in the 3-gallon jug. Assume the initial state to be $(0,0)$ and the goal state to be $(2,0)$. Output the sequence of steps to go from the initial to the goal state.
& & \\
\hline

4. & &
Develop an Expert System for medical diagnosis in Prolog. Given the symptoms, the expert system should predict the disease; given the disease, the expert system should display the symptoms.
& & \\
\hline

5. & &
Write a Prolog Program to append two lists.
& & \\
\hline

6. & &
Write a Prolog Program to reverse a list.
& & \\
\hline

7. & &
Write a Prolog Program to find the $n$th element from a list of integers.
& & \\
\hline

8. & &
Write a Prolog Program to check whether a given element is a member of a list.\newline
\textit{(Note: Use a recursive approach for Prolog programs 9 to 11)}
& & \\
\hline

9. & &
Implement a program to generate the Fibonacci series up to a given number $N$.
& & \\
\hline

10. & &
Implement a program to display the Factorial of a given number $N$.
& & \\
\hline

11. & &
Implement a program to implement the Tower of Hanoi.
& & \\
\hline

12. & &
Implement the Tic-Tac-Toe game using the Minimax algorithm for adversarial search. \textit{(Note: The first move is to be made by humans. The next move should be made by the computer and so on.)}
& & \\
\hline

13. & &
Perform the following:\newline
\quad (i)  Get an intelligent dialogue out of Eliza\newline
\quad (ii) Get an intelligent dialogue out of ChatGPT
& & \\
\hline

14. & &
Parse five English sentences using the Link Grammar Parser.
& & \\
\hline

15. & &
Consider a simple database of books and authors:\newline
\texttt{author(jane, book1).}\newline
\texttt{author(mary, book2).}\newline
\texttt{author(john, book3).}\newline
\texttt{author(jane, book4).}\newline
\quad (i)  Write a Prolog rule to find all books by a given author (including repeated results if the same author wrote multiple books).\newline
\quad (ii) Write a Prolog rule to check if a person is an author. Use cut to prevent unnecessary backtracking.
& & \\
\hline

16. & &
Write a Prolog program to check if a number is a member of a list of integers. Use cut to prevent unnecessary backtracking.
& & \\
\hline

\end{longtable}

\end{document}
